\documentclass[11pt]{article}
\usepackage{amsmath, amssymb, amsthm, geometry, mathtools, booktabs}
\usepackage{bm}
\usepackage{hyperref}
\geometry{margin=1in}

\title{RG-First Conditional Transformer Flow for 2D $\phi^4$: Code-Level Mathematical Description and Design Study}
\author{Kostas Orginos}
\date{\today}

\begin{document}
\maketitle

\section{Purpose}

This note documents the model currently implemented in
\texttt{jaxqft/models/phi4\_rg\_cond\_flow.py}. The old multigrid RealNVP
model in \texttt{jaxqft/models/phi4\_mg.py} is left unchanged. The new model
uses:
\begin{itemize}
\item a recursive $2\times 2$ RG split,
\item conditional affine couplings acting only on the three local fine modes,
\item a local transformer that parameterizes the coupling functions,
\item a learned terminal RealNVP on the last $2\times 2$ block.
\end{itemize}

The note also summarizes the empirical ablation study used to tune the new
design against the old MG-flow baseline at
\[
L=16,\qquad \lambda=2.4,\qquad m^2=-0.4.
\]

\section{Lattice and Prior}

Let the fine lattice be $\Lambda_0$ with shape $L\times L$, where
$L=2^K$. The model is defined for a batch of scalar fields
\[
\phi \in \mathbb{R}^{B\times L\times L}.
\]

The latent variable $z$ is sampled sitewise from the standard normal
distribution,
\begin{equation}
p_Z(z)=\prod_{x\in\Lambda_0}\frac{1}{\sqrt{2\pi}}
\exp\!\left(-\frac{z_x^2}{2}\right).
\end{equation}

The code samples this prior directly on the fine lattice. Because the
``average'' RG split is orthogonal and the ``select'' RG split is a blockwise
permutation, this is equivalent to independent standard normals for all latent
fine modes and for the terminal $2\times 2$ block variables.

\section{Blockwise RG Split}

At RG level $k$ the lattice has shape
\[
L_k\times L_k,\qquad L_k=L/2^k.
\]
For $k=0,\dots,K-2$, the field is partitioned into $2\times 2$ blocks. For
each block we write the packed site values as
\[
v=(v_{00},v_{01},v_{10},v_{11})\in\mathbb{R}^4.
\]

\subsection{Average Rule}

For \texttt{rg\_type="average"}, the code uses the orthonormal Haar-like basis
\begin{equation}
H=
\begin{pmatrix}
\tfrac12 & \tfrac12 & \tfrac12 & \tfrac12 \\
\tfrac12 & -\tfrac12 & \tfrac12 & -\tfrac12 \\
\tfrac12 & \tfrac12 & -\tfrac12 & -\tfrac12 \\
\tfrac12 & -\tfrac12 & -\tfrac12 & \tfrac12
\end{pmatrix}.
\end{equation}
The block coefficients are
\begin{equation}
(c,\eta_1,\eta_2,\eta_3)=H\,v.
\end{equation}
Here $c$ is the normalized coarse mode, while
$\eta=(\eta_1,\eta_2,\eta_3)\in\mathbb{R}^3$ spans the fluctuation subspace.

The arithmetic block average is
\begin{equation}
\bar v = \frac{1}{4}(v_{00}+v_{01}+v_{10}+v_{11})=\frac12\,c,
\end{equation}
which is exactly the quantity used for conditioning in the implementation.

\subsection{Select Rule}

For \texttt{rg\_type="select"}, the coarse variable is the upper-left site and
the fine modes are the remaining three entries:
\begin{equation}
c=v_{00},\qquad \eta=(v_{01},v_{10},v_{11}).
\end{equation}

\subsection{Inverse Block Reconstruction}

At each block, reconstruction is written as
\begin{equation}
v = R(c,\eta).
\end{equation}
For the average rule this is $v=H^\top(c,\eta)$, and for the select rule it is
the obvious concatenation.

\section{Recursive Generative Structure}

Let $z^{(k)}$ denote the latent field on level $k$. The recursion stops when
$L_k=2$, so the terminal latent is a $2\times 2$ field.

For $k=0,\dots,K-2$, split
\begin{equation}
(z_c^{(k+1)},z_f^{(k)}) = S(z^{(k)}),
\end{equation}
where $z_c^{(k+1)}$ is the coarse field on the next lattice and
$z_f^{(k)}\in\mathbb{R}^{3}$ per block.

The generative map is
\begin{align}
x_c^{(K-1)} &= g_{\mathrm{term}}(z_c^{(K-1)}), \\
x_f^{(k)} &= g_k\!\left(z_f^{(k)} \mid x_c^{(k+1)}\right), \\
x^{(k)} &= R\!\left(x_c^{(k+1)},x_f^{(k)}\right),
\end{align}
with the final sample $x=x^{(0)}$.

The inverse map is
\begin{align}
(x_c^{(k+1)},x_f^{(k)}) &= S(x^{(k)}), \\
z_c^{(k+1)} &= g_{\mathrm{term}}^{-1}(x_c^{(K-1)}) \quad\text{at }L_k=2, \\
z_f^{(k)} &= g_k^{-1}\!\left(x_f^{(k)} \mid x_c^{(k+1)}\right), \\
z^{(k)} &= R\!\left(z_c^{(k+1)},z_f^{(k)}\right).
\end{align}

\section{Conditional Fine-Mode Flow}

At each non-terminal level, the flow acts only on the three fine variables of
each block. For coupling layer $j$, the implementation uses one of the masks
\begin{equation}
m^{(0)}=(1,1,0),\qquad
m^{(1)}=(0,1,1),\qquad
m^{(2)}=(1,0,1),
\end{equation}
cycled as needed.

Write the fine vector at one block as $\eta\in\mathbb{R}^3$. Then
\begin{equation}
\eta_A = m\odot \eta,\qquad \eta_B=(1-m)\odot \eta.
\end{equation}
The coupling is
\begin{equation}
\eta_B' = \eta_B \odot e^{s(\eta_A,c)} + t(\eta_A,c),
\end{equation}
while $\eta_A$ is unchanged. In the code, $s$ is clipped as
\begin{equation}
s \leftarrow s_{\max}\tanh(s),\qquad s_{\max}=\texttt{log\_scale\_clip}.
\end{equation}

The inverse is
\begin{equation}
\eta_B = (\eta_B' - t(\eta_A,c)) \odot e^{-s(\eta_A,c)}.
\end{equation}

The log-Jacobian contribution of one coupling layer is
\begin{equation}
\log \left|\det J_j^{-1}\right|
= -\sum_{\text{blocks}}\sum_{a=1}^{3}(1-m_a)\,s_a.
\end{equation}
The full level log-Jacobian is the sum over all coupling layers.

\section{Local Transformer Parameterization}

The coupling functions $s,t$ are not parameterized by a plain MLP. Instead,
for each coupling layer a local transformer processes a block-grid feature map.

\subsection{Token Features}

At each block location $b$, the token feature is
\begin{equation}
u_b = \bigl[m\odot \eta_b,\ \tilde c_b\bigr]\in\mathbb{R}^4,
\end{equation}
where
\[
\tilde c_b=
\begin{cases}
\tfrac12 c_b,& \text{average rule,} \\
c_b,& \text{select rule.}
\end{cases}
\]
Thus the transformer sees both the masked local fluctuation and the local
coarse-field value.

\subsection{Local Attention}

For attention radius $r$, the neighborhood of block $b$ is
\begin{equation}
\mathcal N_r(b)=\{b+\delta \mid \delta=(\delta_y,\delta_x),\ 
\delta_y,\delta_x\in\{-r,\dots,r\}\},
\end{equation}
with periodic wrapping implemented by lattice rolls.

Each transformer block uses standard multi-head attention with a learned
relative-bias table indexed by offset:
\begin{align}
q_b &= W_Q h_b, \\
k_{b,\delta} &= W_K h_{b+\delta}, \\
v_{b,\delta} &= W_V h_{b+\delta}, \\
\alpha_{b,\delta}^{(h)} &=
\mathrm{softmax}_{\delta\in\mathcal N_r(b)}
\left(
\frac{q_b^{(h)}\cdot k_{b,\delta}^{(h)}}{\sqrt{d_h}}
\beta_{\delta}^{(h)}
\right), \\
\mathrm{Attn}(h)_b &=
W_O \Bigl(\sum_{\delta\in\mathcal N_r(b)}
\alpha_{b,\delta}\,v_{b,\delta}\Bigr).
\end{align}

\subsection{Residual Transformer Block}

Each block applies
\begin{align}
h &\leftarrow h + \mathrm{Attn}(\mathrm{LN}(h)), \\
h &\leftarrow h + \sigma(W_g\,\mathrm{LN}(h)+b_g)\odot
\mathrm{MLP}(\mathrm{LN}(h)),
\end{align}
where the MLP is a two-hidden-layer GELU network and $\sigma$ is the logistic
sigmoid.

\subsection{Coupling Output}

After the final transformer block, a small MLP head produces
\begin{equation}
(s,t)=\mathrm{Head}(h)\in\mathbb{R}^6,
\end{equation}
which is split into two three-vectors.

\subsection{$Z_2$-Symmetric Option}

The implementation exposes a parity mode
\[
\texttt{parity}\in\{\texttt{none},\texttt{sym}\}.
\]
For \texttt{parity="none"}, the transformer output is used directly.

For \texttt{parity="sym"}, the code enforces the old-map-style symmetry
\[
z\mapsto -z,\qquad \phi\mapsto -\phi
\]
by requiring the scale channel to be even and the shift channel to be odd in
the conditioning input. If the local transformer head evaluated on token field
$u$ gives
\[
\mathrm{Head}(u)=\bigl(\hat s(u),\hat t(u)\bigr),
\]
then the actual coupling functions are
\begin{align}
s(u) &= \frac12\bigl(\hat s(u)+\hat s(-u)\bigr), \\
t(u) &= \frac12\bigl(\hat t(u)-\hat t(-u)\bigr).
\end{align}
Therefore
\[
s(-u)=s(u),\qquad t(-u)=-t(u),
\]
which implies the desired oddness of the full bijection.

\section{Terminal $2\times 2$ Flow}

The recursion stops at lattice size $2\times 2$. The last coarse block is not
left under a fixed Gaussian anymore. Instead, the code applies a standard
unconditional RealNVP
\begin{equation}
g_{\mathrm{term}}:\mathbb{R}^4\to\mathbb{R}^4
\end{equation}
to the packed terminal block. This terminal flow uses the same mask pattern
$(1,0,0,1)$ / $(0,1,1,0)$ as the older MG-flow implementation, and it inherits
the same parity mode \texttt{none} or \texttt{sym}.

\section{Exact Log Density}

The inverse map returns $(z,\log |\det J^{-1}|)$. Hence the model log-density
is
\begin{equation}
\log p_\theta(x)
=
\log p_Z(z)
\;+\;
\sum_{k=0}^{K-2}\log \left|\det J_k^{-1}\right|
\;+\;
\log \left|\det J_{\mathrm{term}}^{-1}\right|.
\end{equation}

Training minimizes the usual forward KL surrogate
\begin{equation}
\mathcal L(\theta)
=
\mathbb E_{z\sim p_Z}\Bigl[
\log p_\theta(g_\theta(z)) + S(g_\theta(z))
\Bigr].
\end{equation}

\section{Initialization and the Gradient-Flow Fix}

An important empirical point is that exact-zero output initialization for the
transformer head and residual projections is a bad choice here. With exact zero
initialization, the model is identity-preserving, but the initial gradients of
most transformer-internal weights vanish identically. In the measured code path,
the initial norms of representative gradients such as
\texttt{input/W}, \texttt{qkv/W}, and \texttt{mlp/l1/W} were exactly zero, while
only the final output layer had nonzero gradient.

To fix this, the current implementation uses a small nonzero output
initialization scale
\begin{equation}
\varepsilon_{\mathrm{out}} = 10^{-2}
\end{equation}
for:
\begin{itemize}
\item the transformer head final linear map,
\item the attention output projection,
\item the transformer-block MLP final linear map,
\item the terminal RealNVP final linear maps.
\end{itemize}

This preserves a near-identity initial map while allowing gradients to reach
the transformer internals from the first optimization step.

\section{Empirical Design Study}

All experiments below used
\[
L=16,\qquad \lambda=2.4,\qquad m^2=-0.4,\qquad \text{batch}=8.
\]
The comparison metric is the training objective after $20$ optimization steps,
the centered $\Delta S$ standard deviation on held-out flow samples, and the
steady-state step time on CPU.

\subsection{Untuned New Model}

Before the gradient-flow fix, all transformer variants behaved nearly
identically and underperformed the old MG baseline. This was traced directly to
the exact-zero output initialization discussed above.

\subsection{Corrected New Model}

After enabling nonzero output initialization and using learning rate
$3\times 10^{-4}$, the ablation became meaningful.

\begin{center}
\begin{tabular}{lcccc}
\toprule
Model & Loss & $\sigma(\Delta S-\langle \Delta S\rangle)$ & Step [s] & Comment \\
\midrule
Old MG, width 128 & 24.42 & 21.76 & 0.0216 & tuned baseline \\
New, local only, $r{=}0$, 1 blk, 1 head & 85.46 & 30.75 & 0.0333 & too weak \\
New, $r{=}1$, 1 blk, 1 head & 79.01 & 28.86 & 0.0374 & neighbor coupling helps \\
New, $r{=}1$, 1 blk, 4 heads & 79.34 & 28.97 & 0.0383 & heads alone do little \\
New, $r{=}1$, 2 blks, 4 heads & 24.24 & 14.73 & 0.0662 & best tested quality \\
New, select rule, $r{=}1$, 1 blk, 1 head & 93.14 & 33.56 & 0.0376 & worse blocking choice \\
\bottomrule
\end{tabular}
\end{center}

Additional speed/quality sweeps around the best architecture gave:
\begin{itemize}
\item width $32$ or $48$ degrades strongly,
\item reducing the number of couplings from $4$ to $3$ hurts quality,
\item one head and four heads perform similarly in quality once gradients flow,
but four heads were slightly faster on the tested CPU run,
\item the dominant qualitative improvement came from using two transformer
blocks rather than one.
\end{itemize}

\subsection{Effect of Enforcing the $Z_2$ Symmetry}

Using the same tuned architecture and changing only the parity mode gives:

\begin{center}
\begin{tabular}{lcccc}
\toprule
Model & Loss & $\sigma(\Delta S-\langle \Delta S\rangle)$ & ESS proxy & Step [s] \\
\midrule
New tuned, \texttt{parity=none} & 24.24 & 14.73 & 0.00397 & 0.0668 \\
New tuned, \texttt{parity=sym} & 3.42 & 10.60 & 0.00916 & 0.0922 \\
\bottomrule
\end{tabular}
\end{center}

Thus, in the tested short-budget regime, enforcing the symmetry improves both
the loss and the centered reweighting-width proxy, at the cost of a moderate
increase in per-step runtime.

\section{Recommended Design}

Among the tested new-model variants, the best design was:
\begin{itemize}
\item \texttt{rg\_type = average},
\item \texttt{width = 64},
\item \texttt{n\_couplings = 4},
\item \texttt{transformer\_layers = 2},
\item \texttt{n\_heads = 4},
\item \texttt{attn\_radius = 1},
\item \texttt{terminal\_n\_layers = 2},
\item \texttt{output\_init\_scale = 1e-2},
\item \texttt{parity = sym},
\item \texttt{lr = 3e-4}.
\end{itemize}

This configuration improved markedly over the untuned new model and improved on
the tuned old MG-flow baseline in the centered $\Delta S$ metric under the
short test budget. However, it remained substantially slower per training step
on CPU.

\section{Interpretation}

The main lesson from the experiments is:
\begin{enumerate}
\item \textbf{Neighbor coupling matters}, but only after the transformer is
allowed to learn.
\item \textbf{Two transformer blocks matter more than multi-head count.}
\item \textbf{Average blocking is preferable to select blocking.}
\item \textbf{Enforcing the $Z_2$ symmetry helps materially} for this $\phi^4$
task in the tested regime.
\item \textbf{The dominant failure mode of the first implementation was not the
RG factorization itself but the exact-zero output initialization, which froze
the attention stack.}
\end{enumerate}

\section{Open Directions}

The following design directions remain plausible:
\begin{itemize}
\item explicit conditioning on neighboring coarse fields and neighboring masked
fine modes at multiple radii,
\item coupling-order variants that update coarse and fine modes more
symmetrically,
\item a cheaper conditional network than attention if wall-clock speed is the
dominant objective,
\item using the new model as a higher-quality but more expensive proposal and
the old MG flow as a fast baseline.
\end{itemize}

\end{document}
